\documentclass[11pt,letterpaper]{article}

% Enable the inclusion of EPS-format graphics.
\usepackage{epsfig}
% Used to handle line spacing. Must be loaded before hyperref.
\usepackage{setspace}
% Clickable links.
\usepackage[bookmarks,colorlinks=true,linkcolor=blue,urlcolor=blue]{hyperref}
% Code listings via \lstlisting.
\usepackage{listings}
% Multirow table cell spanning.
\usepackage{multirow}
% Zero paragraph indent and non-zero paragraph separation.
\usepackage{parskip}
% Tables with flexible-width columns.
\usepackage{tabulary}
% Diagrams and stuff.
\usepackage{tikz}

% Set document fonts: Times as the default Roman font, Helvetica for
% sans-serif, Courier for teletype.
\renewcommand{\rmdefault}{ptm}
\renewcommand{\sfdefault}{phv}
\renewcommand{\ttdefault}{pcr}

% I want captions to have slightly smaller text and a bold label.
\usepackage[font=small,labelfont=bf]{caption}

% Customize formatting for code listings.
\lstset{
    basicstyle={\small\ttfamily}
}

% LaTeX' default line spacing is a little tight for my liking.
\setstretch{1.15}

\begin{document}

\title{\includegraphics[width=0.75\textwidth]{logo.eps}\\
Reference Implementations Explained}
\author{Samuel Coleman}
\date{\today}

% Turn off index link generation for the unnumbered title page so we don't get
% a warning about a missing page.
\hypersetup{pageanchor=false}
\begin{titlepage}

\maketitle
% Ditch the page styling (including page number) for the title page.
\thispagestyle{empty}

\end{titlepage}
% Having generated the page, we can turn index link generation back on.
\hypersetup{pageanchor=true}

\tableofcontents

\newpage

\section{Introduction}
\label{introduction}

To demonstrate the validity of the Garrit architecture, reference
implementations are provided covering all capabilities. These may form part of
a production environment, a basis for further component development, or simply
be used to demonstrate a functioning example of the system.

\section{Submitter and Reporter: Subport}
\label{subport}

\textbf{Source repository: \url{https://github.com/Garrit/subport}}

Subport provides a web interface through which individuals can submit source
code for evaluation against a problem. In turn, it pushes the submission into
the negotiator. After the submission has been evaluated, it stores the judging
results to disk, and uses this data to provide web interface which shows a list
of all submissions. The user can ``drill down'' into each submission to view
the input source code and the output of the submission for each problem case.

\section{Executor}
\label{executor}

\textbf{Source repository: \url{https://github.com/Garrit/executor}}

The reference executor runs submitted source code in a contained environment by
leveraging Linux containers (LXC). Based on the problem selected for a
submission, the executor feeds input into the program and record its output and
the time consumed to generate said output. The executor also limits execution
time so that programs cannot ``run away'' and monopolize system resources.

\section{Judge}
\label{judge}

\textbf{Source repository: \url{https://github.com/Garrit/judge}}

For each execution case of an execution, the reference judge performs a
byte-level comparison with reference output for that problem case. If the
outputs match, the submission is scored as 100\%; otherwise, it is scored as
0\%.

\section{Negotiator}
\label{negotiator}

\textbf{Source repository: \url{https://github.com/Garrit/negotiator}}

The negotiator provides messaging services between all other components. The
reference negotiator operates somewhat like a dumb pipe: submissions enter,
pass from one component to another, and leave. It does correctly receive error
indications from components and passes those errors along to the reporter, but
it lacks retrial or failover abilities. It serves well in a demonstrative
capacity due to its simplicity.

\end{document}